The peg moves on a circle at a constant speed. Its shadow is its projection onto a line: a harmonic oscillation with the radius of the circle as amplitude and the period of one turn.
\begin{abcliste}
\abc The amplitude is $\result{A = r = \rpO}$. One turn takes
\begin{align}
 T &= \TF = \frac{1}{\n} = \T \approx \result{\TP{0}{3}}
\end{align}

\abc
\begin{align}
 \omega &= \omF = 2\pi\times\n = \om \approx \result{\omP{0}{3}}
\end{align}

\abc A body on a spring oscillates with $\omega = \sqrt{k/m}$. For the same $\omega$ as the shadow:
\begin{align}
 k &= \kDF = \m\times\left(\om\right)^2 = \kD \approx \result{\kDP{0}{3}}
\end{align}

\abc The shadow is fastest when it passes the middle of its path, where the peg moves parallel to the wall. Then it moves as fast as the peg itself:
\begin{align}
 \hat v &= \vmF = \om\times\rp = \vm \approx \result{\vmP{0}{2}}
\end{align}
At that moment the peg is $\result{\text{closest to or farthest from the wall}}$: on the line through the axis perpendicular to the wall.
\end{abcliste}
